\begin{table}[htbp]
\centering\small
\caption{Primary performance of the common-model development procedure and matched controls.}\label{tab:model}
\begin{tabular}{lrrr}\toprule
 & W0 & W12 & W24 \\ \midrule
Patients & 152 & 89 & 75 \\ 
MAE (95\% CI) & $3.19$ ($2.75$, $3.70$) & $1.30$ ($0.96$, $1.71$) & $1.19$ ($0.89$, $1.53$) \\ 
BSA-only area MAE & 3.12 & 1.18 & 1.02 \\ 
Three-marker-only MAE & 7.63 & 7.57 & 8.69 \\ 
\shortstack[l]{Primary minus BSA-only area\\MAE difference (95\% CI)} & $0.07$ ($-0.03$, $0.17$) & $0.12$ ($0.04$, $0.20$) & $0.17$ ($0.11$, $0.24$) \\ 
RMSE & 4.45 & 2.29 & 1.87 \\ 
Mean signed error & $-0.59$ ($-1.27$, $0.07$) & $0.82$ ($0.43$, $1.29$) & $0.84$ ($0.50$, $1.25$) \\ 
Spearman correlation & 0.918 & 0.927 & 0.904 \\ 
Concordance correlation & 0.910 & 0.900 & 0.869 \\ 
Within two points, \% & 41.8 & 84.3 & 82.0 \\ 
Severity-band agreement, \% & 83.2 & 88.2 & 92.7 \\ 
\bottomrule\end{tabular}
\par\small\raggedright W0/W12 use patient-grouped nested validation; W24 is same-cohort temporal application. Structure and penalty selection are internal to each training fold; earlier marker-panel selection is not repeated. Matched controls use the same patients and outer folds. BSA-only area uses the zero-intercept area structure; three-marker-only uses CCL17, CCL26 and IL-5 without BSA. Positive paired MAE differences mean greater primary-procedure error. Metrics are averaged after scoring each repetition. Intervals are 95\% patient-bootstrap intervals conditional on saved predictions; selection and fitting are not repeated. MAE, mean absolute error; RMSE, root mean squared error. Positive signed error indicates overestimation. Severity-band agreement uses four EASI bands for observed and predicted values: $\leq7$, $>7$ to $\leq21$, $>21$ to $\leq50$, and $>50$. The single final formula is not independently validated.\end{table}
